%!TEX root = ../main.tex

\chapter{結論}
本研究では，連続値を扱い高精度な解が得られる線形計画モデルと，期の数の増大に対して計算負荷の伸びが安定している動的計画モデルを構築した．そして，高知県内の工場を対象とした現実的な設定に基づく計算例を通して，各モデルの特徴について以下の点が明らかとなった．
\begin{itemize}
    \item 線形計画モデルおよび動的計画モデルの双方が，売電単価が高い時間帯での売電や計画期間末の蓄電池の使い切りなど，電力運用の最適化という目的に適した合理的な挙動を示すこと．
    \item 動的計画モデルは，計算時間が期の数に対して線形に増加する特性を持ち，長期間の運用計画作成においてスケーラビリティの面でメリットを有すること．
    \item 一方で動的計画モデルは離散化に伴う誤差が懸念されるが，蓄電池の残量の状態の数を適切に設定することで，実用的な時間内で線形計画モデルと遜色のない運用案を導出可能であること．
    \item 線形計画モデルでは，契約電力料金のような複数の期にまたがる制約条件に対しても柔軟に対応可能であるが，動的計画モデルでは，状態空間の次元が増大し計算が困難になるという課題があること．
\end{itemize}

本研究で構築した2つの数理モデルは，太陽光発電量，消費電力量，売買価格などの入力データを差し替えることで，
対象とした特定の工場に限らず，太陽光発電設備と蓄電池を有する他の工場や商業施設などに対しても広く適用可能である．
また，本モデルは運用計画の策定のみならず，設備の導入計画におけるシミュレータとしての応用も期待できる．
例えば，蓄電池容量をパラメータとして変化させながら本モデルによる運用コスト評価を繰り返し行うことで，
導入コストと運用削減効果のバランスを考慮した最適な蓄電池容量の決定といった，設計フェーズの意思決定支援にも寄与し得る．

なお，本研究では太陽光発電量や電力消費機器での消費電力量の予測値が既知であると仮定したが，
実運用においては予測誤差が生じ得る．したがって，発電量や需要の予測値の不確実性を考慮した頑強なモデルを構築することが今後の課題として挙げられる．